% Section 2.5, from lectures/L02/README.md: what you should be able to do, questions to test
% yourself, and the next lecture.

\section{Review}
\label{c2:sec:review}

\needspace{6\baselineskip}
\subsection*{What you should be able to do now}
\begin{itemize}
  \item Say what each of \code{DDRx}, \code{PORTx} and \code{PINx} does, and give the four states
    of a pin.
  \item Configure a pin as an input with its pull-up on, and say what it reads with nothing
    connected.
  \item Explain why writing a one to \code{PINx} toggles an output, and why writing a zero does
    nothing.
  \item Translate an Arduino pin number into a port, a bit, a mask and three register addresses.
  \item Say exactly what \code{rcall} pushes, in what order, and what \code{ret} leaves behind.
  \item Say which registers a subroutine may clobber, and what breaks when it clobbers the others.
  \item Write a driver that reads, modifies and writes a port register rather than assigning to
    it, and say what goes wrong when it assigns.
  \item Predict how a driver's cost varies with the pin it was given, and check the prediction.
\end{itemize}

\needspace{6\baselineskip}
\subsection*{Questions to test yourself}
\begin{enumerate}
  \item \code{DDRx} bit is 0 and \code{PORTx} bit is 1. Is that pin an input or an output, and
    what does it read?
  \item You want to turn one LED on without disturbing the seven other bits of its port. Write the
    three-instruction sequence, and say why a single \code{sts} will not do.
  \item A button is wired from a pin to ground with the internal pull-up on. \code{btn_pressed}
    should return 1 when it is pressed. What does \code{PINx} read at that moment, and what must
    the driver do about it?
  \item \code{rcall} pushes two bytes. Are they the byte address of the next instruction or its
    word address, and which byte goes at the higher address?
  \item Your subroutine uses \code{r16} as scratch and does not restore it. Nothing in this book
    breaks. Give the situation in which something does.
  \item \code{led_on} costs 59 cycles for Arduino pin 13 and 29 for pin 8. Where do the 30 cycles
    go, and what does that tell you about wiring an LED to a low-numbered pin?
\end{enumerate}

\needspace{6\baselineskip}
\subsection*{Looking ahead}
Chapter~3 is about code that runs when you are not looking:
\begin{itemize}
  \item The interrupt vector table, and why a vector holds a jump rather than a handler.
  \item What the hardware saves when an interrupt fires, which is less than you would like.
  \item Pin change interrupts, and one handler serving eight pins.
  \item The button driver, and the first bug in this book that only appears sometimes.
\end{itemize}
